\section{Stakeholder \& Bloc Positions}

\subsection{Summary of Republican Perspectives}

The Republican position on commercial space is characterized by an aggressive push for deregulation, streamlined licensing, and framing U.S. space dominance as a matter of national security—particularly in response to China’s lunar ambitions. The 2024 Republican Party Platform explicitly calls for creating “a robust Manufacturing Industry in Near Earth Orbit,” sending American astronauts back to the Moon and onward to Mars, and enhancing “partnerships with the rapidly expanding Commercial Space sector.”\footnote{Republican National Committee. (2024). \emph{2024 Republican Party platform}. The American Presidency Project. \url{https://www.presidency.ucsb.edu/documents/2024-republican-party-platform}} The Trump administration issued Executive Order 14335 in August 2025, titled “Enabling Competition in the Commercial Space Industry,” directing federal agencies to “eliminate or expedite” environmental reviews (NEPA) for commercial launch and reentry licenses, reform FAA Part 450 licensing regulations, and establish an individualized “mission authorization” process for novel space activities not covered by existing frameworks.\footnote{Executive Order No. 14335, Enabling Competition in the Commercial Space Industry, 90 Fed. Reg. (August 13, 2025). \url{https://www.govinfo.gov/content/pkg/FR-2025-08-19/pdf/2025-15822.pdf}} Republican lawmakers have consistently framed the space race in terms of strategic competition with China. At a September 2025 Senate Commerce Committee hearing entitled “There’s a Bad Moon on the Rise: Why Congress and NASA Must Thwart China in the Space Race,” Chairman Ted Cruz listed maintaining U.S. presence in low Earth orbit, returning humans to the Moon, and preparing for Mars missions as urgent national priorities.\footnote{Space Foundation. (2025). \emph{U.S. Senate Committee on Commerce, Science \& Transportation hearing: "There's a Bad Moon on the Rise."} \url{https://www.spacefoundation.org/reports/u-s-senate-committee-on-commerce-science-transportation-hearing-theres-a-bad-moon-on-the-rise-why-congress-and-nasa-must-thwart-china-in-the-space-race/}} On colonization specifically, the Republican-sponsored House Commercial Space Act of 2023 sought to streamline the regulatory environment and give operators broad latitude to operate in space, with critics noting that the bill would “give operators free rein to operate as they see fit in space” and would “prevent any meaningful regulations to protect astronomy and the space environment.”\footnote{American Astronomical Society. (2023, November). Legislation on space debris and commercialization. \url{https://aas.org/posts/advocacy/2023/11/legislation-space-debris-and-commercialization}}

\subsection{Summary of Democratic Perspectives}

Democrats generally support commercial space growth but raise significant concerns about environmental sustainability, debris mitigation, equitable access, and governance of future settlements. The Biden administration maintained support for the Artemis program but adopted the goal of a “light-touch” regulatory framework through the proposed Mission Authorization process coordinated by the Office of Space Commerce.\footnote{Office of Space Commerce. (2026, March 24). OSC releases updated Mission Authorization proposal. U.S. Department of Commerce. \url{https://space.commerce.gov/osc-releases-updated-mission-authorization-proposal/}}

\textbf{Key Democratic concerns center on several issues.}

First, orbital debris: with annual satellite launches escalating from 200 in 2013 to over 2,600 in 2023, Democrats have emphasized the need for stronger regulations, fiscal interventions, and multilateral institutions for debris mitigation.\footnote{Organisation for Economic Co-operation and Development. (2024, July). Managing space debris to protect our access to the stars. \url{https://www.oecd.org/en/blogs/2024/07/managing-space-debris-to-protect-our-access-to-the-stars.html}} The FAA introduced new preventative measures in 2023 to reduce the long-term impact of upper-stage launch vehicles on the space environment.\footnote{Stanford Space Law Society. (2025). Who takes out the trash in space? Stanford Law School. \url{https://law.stanford.edu/2025/08/22/who-takes-out-the-trash-in-space/}} Second, the FAA’s dual mandate: some members of Congress have raised concerns about whether the FAA can simultaneously promote the commercial space industry and regulate it for public safety.\footnote{Congressional Research Service. (2024). \emph{Regulation of commercial human spaceflight safety: Overview and issues for Congress} (Report No. R48050). \url{https://www.congress.gov/crs-product/R48050}} Third, governance of private settlements: Democrats raise fundamental questions about whose laws apply in commercial habitats, whether companies can claim territorial control, and how the principle of outer space as “the province of all mankind” intersects with private colonization ambitions, particularly given that the Outer Space Treaty prohibits national appropriation but was drafted before commercial settlements were conceivable.\footnote{United Nations Office for Outer Space Affairs. (n.d.). \emph{Treaty on Principles Governing the Activities of States in the Exploration and Use of Outer Space}. \url{https://www.unoosa.org/oosa/en/ourwork/spacelaw/treaties/introouterspacetreaty.html}} Fourth, equitable access: international frameworks like the UN’s Pact for the Future (2024) identify regulating space traffic, mitigating debris, and ensuring equitable access to space resources as critical for sustainable governance.\footnote{European Journal of Futures Research. (2025). Sustainable space governance as a key to protect future generations. \emph{Springer}. \url{https://link.springer.com/article/10.1186/s40309-025-00254-8}}

\subsection{External Stakeholders and Industry Voices}

\subsubsection{Government Agencies}

\begin{itemize}
  \item \textbf{NASA} has transitioned from a monopoly operator to an anchor customer, relying on public-private partnerships (COTS, Commercial Crew) to reduce costs while maintaining safety oversight. NASA continues its transition of low Earth orbit to private industry. On colonization, NASA supports the Artemis Accords framework and views commercial capabilities as essential for sustained lunar presence.\footnote{National Aeronautics and Space Administration. (n.d.). How NASA's work led to commercial spaceflight revolution.\url{https://www.nasa.gov/humans-in-space/commercial-space/how-nasas-work-led-to-commercial-spaceflight-revolution/}}
  \item \textbf{FAA/AST (Office of Commercial Space Transportation)} licenses and regulates commercial launch and reentry. A record 148 FAA-licensed commercial space operations occurred in FY 2024, up 30\% from the prior year, with projections that this may more than double by 2028.\footnote{Executive Order No. 14335, Enabling Competition in the Commercial Space Industry, 90 Fed. Reg. (August 13, 2025). \url{https://www.akingump.com/en/insights/alerts/white-house-issues-long-awaited-order-to-revamp-commercial-space-regs}} Industry stakeholders have called for doubled funding and expedited hiring authority for FAA/AST, arguing the office is “overwhelmed in executing its core launch and reentry mission.”
  \item \textbf{FCC} manages spectrum allocation for satellite constellations, critical for mega-constellations like Starlink and the emerging competition from China’s Guowang constellation.
  \item Office of Space Commerce (Department of Commerce) is developing the proposed “mission authorization” framework for novel activities not covered by existing regulatory regimes, including on-orbit servicing, in-space manufacturing, active debris removal, and cislunar or lunar surface operations.\footnote{Greenberg Traurig, LLP. (2025, October). Executive Order aims to accelerate commercial space development through deregulation. \url{https://www.gtlaw.com/en/insights/2025/10/executive-order-aims-to-accelerate-commercial-space-development-through-deregulation}}
\end{itemize}

\subsubsection{Industry Companies}

\begin{itemize}
  \item \textbf{SpaceX} is the world’s most active launch services provider. In February 2026, Elon Musk announced SpaceX would prioritize establishing a “self-growing city” on the Moon, achievable in less than 10 years, before resuming Mars colonization work in five to seven years.\footnote{Reuters. (2026, February 9). SpaceX prioritizes lunar 'self-growing city' over Mars project. \url{https://www.reuters.com/science/musk-says-spacex-prioritise-building-self-growing-city-moon-2026-02-08/}} SpaceX’s focus on reusable rockets (Falcon 9, Starship) has dramatically reduced launch costs and enabled the commercial revolution, though critics note a “disconnect between the actual work that SpaceX is doing versus the grandiose way that Musk has talked about future colonization plans.”\footnote{The Verge. (2026). Are Elon Musk's Mars plans finally coming back down to Earth? \url{https://www.theverge.com/science/880672/spacex-elon-musk-moon-mars-ipo-blue-origin}}
  \item \textbf{Blue Origin} focuses on a vision of millions of people living and working in space. Along with SpaceX, Blue Origin is pivotal in advancing the economic feasibility of deep-space missions and developing heavy-lift capabilities.\footnote{Hussain, A., et al. (2024). Towards sustainable horizons: A comprehensive blueprint for Mars colonization. \emph{Heliyon}, \emph{10}(4), e25302. \url{https://pmc.ncbi.nlm.nih.gov/articles/PMC10884476/}}
  \item \textbf{Axiom Space} is developing commercial modules for the ISS and plans for a future private space station, representing the transition from government-owned orbital infrastructure to commercial destinations.
  \item \textbf{Rocket Lab} has emerged as a significant commercial launch operator, contributing to the decrease in launch costs and increased access to space.\footnote{Stanford Space Law Society. (2025). Who takes out the trash in space? Stanford Law School. \url{https://law.stanford.edu/2025/08/22/who-takes-out-the-trash-in-space/}}
\end{itemize}
